\section{Null at the Boundary}
\label{sec:ch11-null-at-the-boundary}
\index{nullability!at a federation boundary|(}%

One signature decides most of what follows:

\begin{graphqlsdl}[fontsize=\footnotesize]
_entities(representations: [_Any!]!): [_Entity]!
\end{graphqlsdl}

A non-null list of nullable elements. Every representation you send is
answered, and any of the answers may be nothing. That is the specification's
design and it is the right one: a subgraph handed twenty-five keys should not
have to fail all of them because it cannot answer
one~\autocite{apollo2026subgraphspec}.

\subsection*{One representation that cannot be answered}

\index{nullability!an error scoped to one representation}%
Three representations to Mosaic, the middle one carrying no category, all three
selecting \code{shippingCost}. The resolver in
section~\ref{sec:ch11-a-field-from-the-other-service} throws for the middle
one, on purpose, because a shipping cost worked out from a category nobody sent
would be a wrong number rather than a missing one:

\begin{minted}[fontsize=\footnotesize]{text}
entities: 5.9, null, 39
error: Unexpected Execution Error at _entities.1.shippingCost
\end{minted}

\code{Money!} cannot be null, so the null propagates up to the \code{\_Entity}
element, which is nullable and stops it there. The neighbours are untouched.
The index in the path is how the client works out which of its representations
failed, and it is the only thing in the response that says so.

\index{errors!Unexpected Execution Error@\code{Unexpected Execution Error}}%
The message is worth a moment. \code{Unexpected Execution Error} is what
HotChocolate says about an exception it was not told how to describe, and the
sentence I wrote into that \code{InvalidOperationException} is in the server
log and nowhere else. That is correct behaviour for a production service, and
the consequence is worth knowing in advance: the client is handed something
that means \enquote{one of your subgraphs threw}, and the text that says which
one and why is in a different process from the one the client is talking to.

\subsection*{Null and null are not the same}

\index{reference resolver!returning null against throwing}%
Now the same batch with a key that does not decode instead of a category that
is missing. Chapter~\ref{ch:the-first-cut-extracting-catalog}'s gate has
asserted this since that chapter and I had never put it beside the case above:
the element is null, and there is no error at all.

Two failures, two identical-looking holes in the data, and one of them is
silent. The difference is what the reference resolver did. Returning null is a
subgraph saying \enquote{not mine}, which the specification treats as an
ordinary answer. Throwing is a subgraph saying \enquote{something went wrong},
which produces an error beside the null. Both are legitimate and you choose
between them every time you write one of these:

\begin{minted}[fontsize=\footnotesize,breaklines]{csharp}
=> ProductKey.TryDecode(id, context, out var productId)
    ? new Product { Id = productId }
    : null;
\end{minted}

Mosaic returns null there because a key it cannot read is genuinely not its
product. If it threw, every router that ever mangled an identifier would fill
your logs on the wrong side of the boundary.

\subsection*{What a client sees, which is worse}

\index{nullability!propagation through a non-null chain}%
Everything above is what the subgraph answered. What the client gets depends
entirely on the nullability of the path the entity hangs from, and that is
where a small local failure stops being small.

The sample has a crate holding a widget key Catalog has never heard of, which
is what a dangling foreign key looks like once the two tables are two services.
Ask for it through the field with no \code{@provides} on it, so the router
really does go and ask:

\begin{minted}[fontsize=\footnotesize,breaklines]{text}
data: null
error: Cannot return null for non-nullable field 'Query.crates.widgetByKey.name'.
path: crates.3.widgetByKey.name
\end{minted}

One crate out of four, and the whole response is \code{data: null}. Follow the
chain outwards from the failure:

\begin{graphqlsdl}[fontsize=\footnotesize]
name: String!   inside   widgetByKey: Widget!   inside   crates: [Crate!]!
\end{graphqlsdl}

Every link is non-null, so there is nowhere for the null to stop and it walks
out of the document.

\index{nullability!the argument for nullable boundary fields}%
This is the measured version of something
chapter~\ref{ch:the-first-cut-extracting-catalog} claimed and did not test. The
doc comment on \code{OrderLineNode.GetProduct} says a dangling identifier
\enquote{surfaces as a null product at the router}. It surfaces as a null
product only where the field is nullable. Where it is not, it surfaces as a
null response, and the three crates that were perfectly fine go with it.

So the rule I would give someone extracting a service, and I did not hold it
myself until this chapter measured it: \textbf{a field that crosses a subgraph
boundary should be nullable unless you can guarantee referential integrity
across a network}. You almost certainly cannot. Mosaic's own
\code{OrderLine.product} is non-null and would take an order query down the same
way, and the only reason it has never happened is that
the seed data is consistent by construction.
Chapter~\ref{ch:strangling-the-monolith} carves four more domains out and will
have to decide this field by field; I would rather it decided with this listing
in front of it.

\index{provides@\code{@provides}!hides a broken reference}%
And notice what section~\ref{sec:ch11-the-promise-provides-makes} would have
done here. Through the \code{@provides} path the same broken reference answered
\enquote{Unknown part} and nobody was any the wiser. The directive that saved a
round trip is also the directive that turned a loud failure into a quiet wrong
answer. Neither behaviour is a bug. Both are the graph doing exactly what its
schema told it to.

\medskip
That is the boundary, in both directions: what the router sends, what a
subgraph does with it, and what comes back when it cannot.
Chapter~\ref{ch:strangling-the-monolith} takes the other five domains out of
\code{Mosaic.Api}, and every one of them will hang fields off a type it does not
own, which is the arrangement this chapter has been taking apart.

\index{nullability!at a federation boundary|)}%
